\section{DAO Governance Mechanisms}


Zoo's DAO operates under non-profit constraints while maintaining decentralized decision-making:

\subsection{Governance Structure}

\begin{enumerate}
\item \textbf{Token Holder Governance}: Quadratic voting (Section 4.1)
\item \textbf{Technical Committee}: 7 elected developers (serve 1-year terms)
\item \textbf{Experience Curation Board}: 11 domain experts (evaluate experience quality)
\item \textbf{Treasury Management Committee}: 5 elected members (oversee DAO spending)
\item \textbf{Executive Director}: Appointed by 501(c)(3) board (ensures legal compliance)
\end{enumerate}

\subsection{Proposal Types and Thresholds}

\begin{table}[h]
\centering
\begin{tabular}{lccc}
\toprule
\textbf{Proposal Type} & \textbf{Quorum} & \textbf{Threshold} & \textbf{Voting Period} \\
\midrule
Protocol parameters & 5\% & 51\% & 7 days \\
Treasury <1B $\KEEPER$ & 3\% & 51\% & 5 days \\
Treasury >1B $\KEEPER$ & 10\% & 66\% & 14 days \\
Experience removal & 2\% & 51\% & 3 days \\
Validator slashing & 5\% & 66\% & 7 days \\
Constitution amendment & 15\% & 75\% & 21 days \\
\bottomrule
\end{tabular}
\caption{DAO proposal requirements by type}
\label{tab:dao_proposals}
\end{table}

\subsection{Quadratic Voting Mechanics}

Zoo implements quadratic voting to balance large and small holders:

\begin{equation}
V_i = \sqrt{S_i}
\end{equation}

where $V_i$ = voting power, $S_i$ = staked $\KEEPER$ tokens

\textbf{Example:}
\begin{itemize}
\item Alice: 100M $\KEEPER$ staked → $\sqrt{100M} = 10,000$ votes
\item Bob: 1B $\KEEPER$ staked (10× more) → $\sqrt{1B} = 31,623$ votes (only 3.16× more)
\end{itemize}

This prevents whale dominance while still rewarding stake.

\textbf{Vote Delegation:}

Users can delegate voting power to experts:
\begin{itemize}
\item Delegation: Transfer voting power (not tokens)
\item Revocable: Instant revocation (no lockup)
\item Public: Delegation visible on-chain (transparency)
\item Split: Can delegate to multiple representatives with percentage weights
\end{itemize}

\subsection{Treasury Management Transparency}

All DAO treasury transactions are:
\begin{itemize}
\item \textbf{On-chain}: Publicly auditable via block explorer
\item \textbf{Proposal-linked}: Each transaction references approved proposal
\item \textbf{Multi-signature}: Requires 3-of-5 Treasury Committee signatures
\item \textbf{Time-locked}: 48-hour delay between approval and execution (allows community intervention)
\end{itemize}

\textbf{Quarterly Reporting:}

Treasury Management Committee publishes quarterly reports:
\begin{itemize}
\item Asset allocation breakdown
\item Quarter-over-quarter spending analysis
\item Grant recipient outcomes
\item Long-term sustainability projections
\end{itemize}

\subsection{Non-Profit Compliance}

Zoo's 501(c)(3) status imposes legal constraints on governance:

\begin{enumerate}
\item \textbf{IRS Oversight}: Annual Form 990 filing (public disclosure)
\item \textbf{Mission Alignment}: All DAO votes must serve charitable purpose (AI democratization)
\item \textbf{No Private Benefit}: Cannot distribute profits to private interests
\item \textbf{Executive Director Veto}: ED can veto DAO decisions that violate 501(c)(3) rules (override requires 80\% supermajority)
\end{enumerate}

This creates \textbf{hybrid decentralization} - maximum autonomy within legal boundaries.

